\part{Self-Overcoming and the Revision of Models}

\chapter{Learning Within a World and Changing the Terms of a World}

Three activities are easily confused under the single name of learning, and the mathematical centre of this book depends on keeping them apart. State estimation is the inference of the current hidden state of the world from observations, given a model. Parameter learning is the adjustment of the quantities that fix how states generate observations and how they evolve, again within a model whose form is held fixed. Structural revision is the alteration of the form itself, so that the variables, the dependencies among them or the kinds of cause admitted are different afterwards. The first two operate inside what may be called an admissible family of models, and the third replaces the family. An agent may improve its estimates indefinitely while retaining an inadequate family, and excellent optimization within that family cannot supply distinctions the family excludes.

Let $\mathcal{M}$ denote the admissible family. Ordinary updating selects within it, producing at each time a posterior over states and parameters that depends on the data seen so far. Structural revision is a change of the family,
\[
\mathcal{M}_{t}\longrightarrow\mathcal{M}_{t+1},
\]
and the question the chapter pursues is what could make such a transition justified. Two kinds of evidence bear on it. The first is a persistent residual. Suppose that over a long run of data the best achievable predictive error within $\mathcal{M}_t$, after all parameter learning, remains above what the noise level of the observations would explain. The excess is then not attributable to ignorance of the current state or to poorly tuned parameters, since both have been optimized away, and it indicates that the family omits some structure the data contain. The second kind of evidence is comparative. If a candidate family $\mathcal{M}'$ is available, its marginal likelihood on the same data can be compared with that of $\mathcal{M}_t$, and the ordinary machinery of model comparison applies.

The comparative route presupposes what the question was meant to explain. A candidate family must already be specified for its evidence to be computed, so that Bayesian model comparison ranks proposals but does not generate them. The residual diagnoses a lack without saying what would fill it. This yields a precise statement of a difficulty that Nietzschean language tends to leave in metaphor, though a narrower one than a general impossibility claim would suggest. Optimization within a specified family does not, by itself, explain how an agent acquires alternatives outside that family. Formal systems can be expanded to include structural alternatives, procedures that propose them, or searches over descriptions, as work on structure learning, nonparametric Bayesian models and program induction does, and such systems do generate new families in a precise sense. The question then concerns what the expanded system presupposes, namely which space of structures is searched, by what proposal mechanism and under what prior, and how its proposals become available to the organism. Relative to a given family, structural novelty is a reorganization of the space in which inference takes place, and an account of the organism's situation, whether encounter with resistance, metaphor, other persons or the inheritance of a language, must say what makes the enlarged space and its proposals available.

Cost enters the decision. Revising a family is not free, since it disrupts the habits, expectations and coordinated action built on the old one. A simple expression of the trade-off is that revision is warranted when the expected improvement in performance over the relevant horizon exceeds the cost of transition,
\[
\E\!\left[\Delta L\right]\,H \;>\; c ,
\]
where $\Delta L$ is the per-period reduction in expected loss, $H$ the horizon and $c$ the transition cost. This inequality is deliberately crude, and its crudeness is instructive. The terms $\Delta L$ and $H$ are themselves estimated from within the current interpretation, so that an organization may undervalue a change whose benefits its own categories cannot register. Cost is also unevenly distributed across agents, and Chapter~11 asks what happens to $c$ when the burden of revision can be shared.

Self-overcoming enters as a philosophical interpretation of certain such reorganizations and not as a synonym for them. The computational notion is value-neutral, since a family can change for better or worse and the criteria of residual and evidence are purely predictive. Nietzsche's evaluative conception of becoming asks which reorganizations express an enhancement of life and which a decline, and this requires a standard that the predictive criteria do not contain. The text preserves the difference between computational revision and evaluative becoming throughout, and Chapter~12 returns to what the standard would have to be.

\chapter{Defensive Coherence and Interpretive Entrapment}

An interpretation can maintain its coherence by reducing its exposure to correction. The mechanisms are familiar from the preceding chapters: selective attention that assigns low precision to disconfirming channels, constrained action that avoids situations where the interpretation is at risk, reinterpretation of counterexamples so that they confirm what they appeared to challenge, and dependence on social environments that reward repetition of the received view. A system organized in this way can become increasingly effective at preserving its interpretation while becoming less capable of handling changes beyond the protected domain. Its reported coherence improves as its adequacy to the wider situation declines.

Nietzsche's analysis of the ascetic ideal supplies an unusually rich case. In the third essay of \emph{On the Genealogy of Morality}, especially in sections 7 and 8, the philosopher's withdrawal from ordinary pleasures and obligations is presented as serving the conditions of intellectual work, so that renunciation can express a powerful commitment to a task. The broader ascetic ideal, however, can turn suffering into an indictment of earthly existence and make detachment a doctrine about life as such. Outwardly similar acts of detachment can accordingly serve very different dispositions: one a discipline that enables activity, the other a defence that forecloses it. The concluding section of the essay presses the point further by observing that human beings would sooner will nothingness than not will. An apparent escape from attachment can preserve an attachment to purity, to salvation or to the identity of one who has transcended ordinary needs. The will has changed its object, and perhaps turned against itself, without ceasing. On the reading developed here, this is a precise instance of defensive coherence: an organization that has secured its stability by making its own renunciation the object it cannot give up.

Dynamical systems supply the main mathematical vocabulary. Consider an interpretive state $x$ evolving under a potential $V$ with noise,
\[
dx=-V'(x)\,dt+\sigma\,dW_t .
\]
Minima of $V$ are attractors, each with a basin of attraction, and the noise occasionally drives the state across the barrier separating one basin from another. For small $\sigma$ the expected time to escape a basin whose barrier has height $\Delta V$ grows as $\exp(2\Delta V/\sigma^{2})$, so that stability in the sense of persistence can be extremely high. The central point is that this measure of stability is silent about adequacy. A deep basin is deep because the barrier around it is high, which says nothing about whether the state at its bottom handles the wider range of situations that the organism may meet. Local stability supplies no sufficient measure of wider adequacy.

Hysteresis and metastability sharpen the distinction between persistence and flexibility. A system exhibits hysteresis when the path by which it moves between states depends on its history, so that returning conditions to an earlier value does not restore the earlier state, and an interpretation that has been reinforced may resist reversal even after the circumstances that justified it have passed. Metastability describes states that persist for long times without being the deepest available, and it captures an interpretation that is held firmly although a better organization exists elsewhere in the space. These notions allow the chapter to state its caution against the opposite error. It would be a mistake to treat instability as intrinsically emancipatory. A system that moves readily among basins may be incapable of sustaining any organization long enough to act, and the capacity to hold an interpretation is as much a condition of life as the capacity to leave it. What distinguishes a healthy organization is not low barrier height but the correspondence between its barriers and the actual costs and risks of the transitions they guard, and a barrier that has become decoupled from those costs is the formal mark of entrapment.

\chapter{Attachment and the Conditions of Revisability}

Attachment theory enters the argument as a question about the conditions under which an organism can tolerate uncertainty and reorganization. In the tradition founded by Bowlby and elaborated empirically by Ainsworth and her colleagues, early relationships with reliable caregivers are understood to support the capacity to explore, to return to a secure base and to regulate distress, and later work extended these ideas to adult relationships \citep{bowlby1969,ainsworth1978}. The chapter uses this tradition in a restricted way. It avoids diagnosing Nietzsche, and it does not treat solitude as an attachment classification. Its subject is the possibility that reliable relationships help sustain the capacity to explore, to receive correction and to undergo changes that temporarily disrupt self-understanding.

Two questions that are easily conflated must be separated at the outset. Nietzsche's criticism of the ego concerns whether there is a simple, enduring subject behind thought and action, and his positive proposal in \emph{Beyond Good and Evil}, sections 12 and 19, understands the soul as a social structure of drives and willing as a relation of commanding and obeying within us. The attachment-theoretical question concerns how a person develops the capacity to explore, to tolerate vulnerability and to sustain relationships. The two are logically independent. Rejecting a substantial ego does not establish emotional independence, and needing relationships does not establish a substantial ego. A self understood as an organization of drives could depend profoundly on relationships for its formation and stability, and nothing in the multiplicity thesis forbids this.

This sharpens the earlier discussion of the ego. If an identity coordinates a person's expectations and actions, as Chapter~3 argued, then challenging it threatens more than a proposition. A criticism of a belief can feel like an injury to the person because the belief helps hold the identity together, and another person's independence can feel threatening where a sense of worth depends on being indispensable. In the terms of Chapter~9, the transition cost $c$ of revising an identity is high exactly because so much coordinated activity depends on it. Relationships and institutions may distribute the burden. If the cost of revision is shared with those who can offer reassurance, correction and continuity during the period of disruption, the effective cost falls,
\[
c_{\mathrm{eff}}=c-s,
\]
where $s\ge 0$ denotes the support available, and changes that would fail the inequality of Chapter~9 for an isolated person may pass it for a supported one. The expression is a bookkeeping device with no empirical content of its own, but it makes a substantive claim visible: autonomy can be understood partly through the availability of support for transformation, and an account that celebrates independence while omitting the supports that produced it describes only half of the conditions of revisability.

This yields the strongest question that can be put to Nietzsche in this area, namely that his account of self-overcoming may underdescribe the relational conditions that make self-overcoming possible. Someone's capacity to endure solitude, withstand disagreement or revise an identity might have been cultivated through reliable relationships, and a portrayal of the resulting independence as self-generated could overlook the support that helped produce it. The question is legitimate without first identifying solitude with avoidance, and it should not be mistaken for the claim that Nietzsche himself was insecurely attached, which no textual evidence can establish.

The paper \emph{Beyond-Man and Relational Agency: A Reading of Zarathustra through Attachment Theory} \citep{ferrari} argues that Zarathustra's solitary self-overcoming resembles insecure-avoidant agency, in that independence becomes defensive self-sufficiency, relationships supply little corrective feedback and vulnerability receives insufficient acknowledgment. It presents its reading as structural and interpretive and disclaims any diagnosis of Nietzsche or his characters, a restraint that the present chapter shares, and it draws its account of avoidant attachment from Bowlby and from Mikulincer and Shaver \citep{mikulincer2007}. Its central question deserves retention: whether apparent autonomy conceals an inability to depend on others. The verdict is less secure than the question. Describing Zarathustra's agency as absolute, unmodulated and largely unresponsive to others requires considerably more textual argument than the paper provides, and the next part, in Chapter~14, tests it against the evidence of the narrative. The illuminating distinction that the paper's question suggests is between independence that permits relationship and independence that must continually defend itself against relationship. Applied to the will to power, it asks whether self-mastery increases a person's capacity to encounter and be changed by others, or whether maintaining the identity of a self-sufficient person becomes an overriding constraint.

\chapter{Value Creation and the Problem of the Objective}

The deepest difficulty in the comparison concerns what evaluates changes in what is valued. In the formalism of Chapter~5, preferences are encoded in the distribution $p(o)$ over preferred outcomes, which appears in the expected free energy of policies and which is part of the generative model and not derived from it. A model that optimizes action relative to specified preferences has not yet explained the authority, formation or transformation of those preferences. If the preferences used to evaluate policies can themselves change, one must ask what evaluates those changes.

One response is to place the preferences in a higher level of a hierarchical model, so that they are themselves subject to learning. A prior over preferences $C_0$ is then updated in light of experience by some operator $R_0$ that depends on a still higher set of commitments $C_1$, and so on. One way of closing such a hierarchy ends in a level whose commitments are held fixed, and an account built in this way can represent changes in preference only by letting an unrevised evaluator carry the authority for everything below it. This is not the only possible architecture. An implemented system must have specified dynamics, but specified dynamics need not amount to an unchanged evaluator that ranks every transformation, and mutually modifying processes, in which no level is exempt from revision by the others, and path-dependent changes are not excluded by finitude. What the hierarchical case does bring out is the difference between explanation and justification. A formal account of how preferences change does not by itself justify those preferences or confer normative authority on the process that changes them, and the question of authority has to be asked of whichever architecture is adopted.

Nietzsche's value creation should be approached through this problem while resisting its reduction to adjusting a reward function. At least three kinds of change are involved and the formalism treats them unevenly. A change in desired outcomes alters which states the agent seeks and is the case that the hierarchy above represents most easily. A change in the relative authority of concerns alters which of several objectives governs conduct, which in the notation of Chapter~4 is a change in the weights $w_i$ and not in the objectives $F_i$ themselves. A change in what counts as an intelligible goal alters the space of possible objectives, so that outcomes that could not previously be wanted become available to want, and this is a structural revision of the kind described in Chapter~9 applied to the objective itself. The third kind is the one that the scalar picture handles worst, because a scalar objective is defined over a given space of outcomes and cannot represent the enlargement of that space as a movement within it.

The ascetic ideal again furnishes a test. When the will turns from worldly objects to willing nothingness, as the \emph{Genealogy} describes, it is difficult to say that a preference has been reweighted, since what has altered is the way the will conceives of its own aims. The case suggests that will to power in Nietzsche's use is better understood as the organizing activity within which preferences are formed than as a preference of especially high weight. If so, it does not occupy the role that $p(o)$ occupies in active inference, and the temptation to identify them rests on a category mistake. The appropriate comparison is between the organizing activity and the process by which a model's preferences, attentional weights and structure are jointly settled, which the standard formalism leaves outside its scope.

Whether formal accounts can represent these transitions without silently installing an unchanged higher-order evaluator remains the open question of the book. The argument of the remaining parts does not claim to resolve it. It claims that the question is correctly posed here, that the formalism supplies exact tools for describing the preservation of an organization and the revision of a model within it, and that the difference between preserving an interpretation that makes life manageable and transforming it so that new forms of life become possible depends on an account of value that the formalism, as ordinarily stated, does not itself supply.
